\subsection*{A determinant estimate for the short-shift rows}

Write
\[
 e=S_0=\epsilon_0=k_0=\frac{30869}{5000000},\quad
 a_0=\frac6{125},\quad \nu=\frac1{64},\quad \xi=\frac1{10000},
\]
\[
 q_-=\frac23-e,\quad q_+=\frac23+2e,\quad
 \tau_0=1-\frac32e .
\AutoTag{eq:old:5_10}
\]
Thus \(q_-=9907393/15000000\), \(q_+=5092607/7500000\), and
\(\tau_0=9907393/10000000\). The high-source inequalities give
\[
 m>\frac{Y^{\tau_0}}{\sqrt3},\qquad
 Y^{q_-}<M<\kappa_Y^2Y^{q_+},\qquad \kappa_Y=1+Y^{-1}.
\AutoTag{eq:old:5_11}
\]
The first inequality follows from \eqref{eq:old:2_8}. The lower bound for \(M\) follows
from failure of the \(O_M\) predicate in the first-match partition
\eqref{eq:old:3_1}. Its upper bound follows from failure of \(O_R\) and
\eqref{eq:old:2_7}.

We first give the determinant estimate for each fixed shift and
coefficient box. Suppose \(r_0,s_0\ge1\),
\(K_{\rm det}=(r_0+1)(s_0+1)\), and
\(K_{\rm det}\) source points lie in one cell. In both applications below, scale
the two projective coordinates by the dyadic box centers. Then a determinant
of the \(K_{\rm det}\) monomial rows has the form
\[
 |\Delta_{\rm det}|\le C_{r_0,s_0}(H_1H_2)^{r_0K_{\rm det}/2}(EF)^{s_0K_{\rm det}/2}
 \max_{\substack{n_\ell\ge0\\\sum_\ell n_\ell=K_{\rm det}}}
 \vartheta^{\sum_\ell \ell n_\ell}
 \lambda^{\sum_\ell n_\ell(n_\ell-1)/2}.
\AutoTag{eq:old:5_12}
\]
Here \(E,F\) are coefficient scales and \(H_1,H_2\) are corner scales.
The cell width in normalized coordinates is
\(\lambda\), and \(\vartheta\) bounds the thickness of the source curve.
In the \(u\)-chart, the remaining functions in a level-\(\ell\) minor are
\(x^{i-2j}\), \(x\in[1/2,2]\). Here
\(\vartheta\) bounds \(\theta\) in \((1-\theta)^j\). In the \(t\)-chart,
put
\[
 \sigma_t=H_B/H_C,\qquad x=(B/C)/\sigma_t,\qquad
 \kappa_t=\sigma_t^2E/F,\qquad y=(E/F)t/(\gamma C^2).
\]
Then \(x\in[1/2,2]\), \(1/16<\kappa_t<8\), and the normalized
coefficient ratio is
\[
 (E/F)(\delta/\gamma)=\kappa_t x^2+y.
\]
The level-\(\ell\) functions are therefore
\[
 x^i(\kappa_t x^2+y)^j
 =\sum_{\ell=0}^j\binom j\ell\kappa_t^{j-\ell}y^\ell
     x^{i+2j-2\ell}.
\]
Fix the determinant degree. The derivatives in \(x\) of these
finitely many monomials are then uniformly bounded, since \(x\in[1/2,2]\)
and \(\kappa_t\in[1/16,8]\). The thickness is the normalized term \(y\).
It is comparable to \(\Xi_t^{-1}\) in the \(3t\) application below.
This normalization gives the same level and cell-width gains in \eqref{eq:old:5_12}.
To verify \eqref{eq:old:5_12}, expand each monomial column as a sum of
level columns and use determinant multilinearity in the columns. In each
resulting determinant, group columns with the same level and apply Laplace
expansion over row subsets of the corresponding sizes. In a level-\(\ell\)
minor, factor the \(\ell\)-th power of the thickness coordinate from each
evaluation row. The resulting thickness exponent is \(\sum_\ell \ell n_\ell\).
If two evaluation points or two powers agree,
its remaining determinant is zero. Otherwise take \(n\) evaluation points
in an interval of length \(\lambda\). Take \(n\) functions with bounded
derivatives through order \(n-1\). Successive divided differences factor
out the Vandermonde product. The remaining determinant is at most a
constant times the derivative bounds. This gives
\(O(\lambda^{n(n-1)/2})\). Thus this group contributes
\(C_{r_0,s_0}\lambda^{n_\ell(n_\ell-1)/2}\). Here all derivatives of
the finitely many powers are bounded on the fixed compact interval. Row
and column scaling gives the first factor in \eqref{eq:old:5_12}.

Put \(b_{\rm cell}=\log_Y(1/\lambda)\) and
\(\mu_{\rm th}=\log_Y(1/\vartheta)\). In the source boxes these exact
quantities equal their limiting exponents up to \(O(1/\log Y)\). Then
\[
\sum_{\ell\ge0}\left(\mu_{\rm th}\ell n_\ell+
 \frac {b_{\rm cell}}2 n_\ell(n_\ell-1)\right)
\ge \frac{2\sqrt2}{3}\sqrt{b_{\rm cell}\mu_{\rm th}}\,K_{\rm det}^{3/2}
-\left(\mu_{\rm th}+\frac {b_{\rm cell}}2\right)K_{\rm det}.
\AutoTag{eq:old:5_13}
\]
For every \(\Lambda>0\), termwise completion of the square gives
\[
\sum_\ell\left(\mu_{\rm th}\ell n_\ell+\frac {b_{\rm cell}}2n_\ell^2\right)
\ge \Lambda K_{\rm det}-\frac1{2b_{\rm cell}}\sum_{\ell\ge0}(\Lambda-\mu_{\rm th}\ell)_+^2
\ge \Lambda K_{\rm det}-\frac{\Lambda^2}{2b_{\rm cell}}-\frac{\Lambda^3}{6b_{\rm cell}\mu_{\rm th}}.
\]
Taking \(\Lambda=(2b_{\rm cell}\mu_{\rm th}K_{\rm det})^{1/2}\) gives the leading
term in \eqref{eq:old:5_13}. The square-completion loss is
\(\Lambda^2/(2b_{\rm cell})=\mu_{\rm th}K_{\rm det}\). Replacing \(n_\ell^2\) by
\(n_\ell(n_\ell-1)\) costs a separate \((b_{\rm cell}/2)K_{\rm det}\).
For \(r_0=\lceil s_0 q/p\rceil\) or \(r_0=\lfloor s_0 q/p\rfloor\),
\[
\frac{r_0p+s_0q}{2}\le s_0q+\frac p2,\qquad
K_{\rm det}^{1/2}\ge (r_0s_0)^{1/2}.
\]
Here \(p\) and \(q\) are the logarithmic scale exponents of the corner
product and coefficient product in the chosen chart. The \(3u\) and \(3t\)
applications below specify their normalizations. Hence a fixed
positive margin in
\[
b_{\rm cell}\mu_{\rm th}>\frac98pq
\AutoTag{eq:old:5_14}
\]
makes every integerized determinant have absolute value less than \(1\).
Here \(p,q,\mu_{\rm th}\) range over fixed compact intervals in both
applications. The strict margin is uniform over the full source boxes.
Fix that margin and those compact intervals, and then choose \(s_0\).
The degree-rounding and lower-order terms are then smaller than the margin.
This one fixed degree works for every sufficiently large \(Y\).
The integerized matrix has rank below \(K_{\rm det}\). Let \(r<K_{\rm det}\)
be the largest order of a nonzero minor. If \(r=0\), a coordinate vector
is a kernel vector. If \(r>0\), choose a maximal nonzero \(r\times r\)
minor and one additional column. Cramer's rule gives a nonzero integral
kernel vector. Its coordinates are minors of order at most \(r\). The selected rows
span the row space, so this vector annihilates every row. Thus the kernel
vector has polynomially bounded height and gives a nonzero integral
auxiliary polynomial.
If a cell contains fewer than \(K_{\rm det}\) points, those points are counted
directly. The degrees are fixed. The coefficient height of the
polynomial is a fixed power of \(Y\). This holds because it is a minor of
a matrix whose entries have fixed polynomial height.

Heath--Brown's Lemma~2~\cite[Lemma~2, p.~8]{HB2012}, stated in
Appendix~\ref{app:sources}, has no size restriction on the second pair.
In \(3u\) and \(3t\), the
primitive coefficient pair is the first pair and the primitive corner
pair is the second. For an irreducible factor, let \(a_c\) be its
degree in the corner pair. Let \(b_q\) be its degree in the coefficient
pair. In the theorem's notation this is \((a,b)=(b_q,a_c)\). Thus a
factor with \(a_c\ge2\) has second-pair degree at least two.
Any application of Heath--Brown below has \(b_q\ge1\). So both
bidegrees are positive. The theorem's exponent is \(2/a_c\le1\). Here
\(T_1+1\) and the factor height are bounded by fixed
powers of \(Y\). Choose the theorem's \(\varepsilon\) after those height
exponents. This gives \(O_\eta(Y^\eta(T_1+1))\).

Let a bounded-degree integer polynomial have height \(Y^C\). Then each
primitive integer factor has height \(Y^{O_C(1)}\). For the proof, apply
a fixed Kronecker substitution with base larger than the total degree.
Then use Gauss's lemma and the univariate Landau--Mignotte factor bound. The
substitution has no collisions among the bounded-degree monomials.

\subsection*{The short-\(u\) row \(3u\)}

In the \(u\)-chart put \(\alpha=ac\), \(\beta=bd\). The source identities
\eqref{eq:old:2_2} give
\[
 \beta D^2-\alpha A^2=u,\qquad
 \alpha A^2=(1-\theta)\beta D^2,\qquad
 \theta=\frac{m+n+1}{(m+1)(n+1)}<\frac2{m+1}.
\AutoTag{eq:old:5_15}
\]
The residual factors are pairwise coprime, and so are the corner
factors. Hence
\[
 \gcd(\alpha,\beta)=\gcd(A,D)=1.
\AutoTag{eq:old:5_16}
\]
Also
\[
 \frac{um}{3}\le\alpha A^2<um,\qquad
 \frac{um}{3}\le\beta D^2\le u(m+1),
\AutoTag{eq:old:5_17}
\]
because \(n/(m+n+1)\ge1/3\). Multiply these two scale identities
and use \(\alpha\beta=M\). This gives
\[
 AD\ge \frac{um}{3\sqrt M}
 \gg uY^{2/3-5e/2}.
\AutoTag{eq:old:5_18}
\]
The implied constant is absolute for sufficiently large \(Y\). The unit
conditions needed below are also source consequences. If a prime divides
\(A\), it divides \(m,n\). So it divides neither \(b\mid m+1\) nor
\(d\mid n+1\). Hence \(\gcd(\beta,A)=1\). Also \(uBC=m+n+1\) and
\(\gcd(A,BC)=1\), so \(\gcd(u,A)=1\). Therefore
\(\gcd(\beta u,A)=1\). Every corner is at most \(Y+1\), and
\(\alpha,\beta\le u(Y+1)\).

Fix a positive shift \(u\) and a dyadic sector. The counted tuple
is \((\alpha,A,D)\), with \(A,D\) in that sector. The conic determines
the second coefficient \(\beta\), subject to the integrality filter.
Fix \(u\le Y^{a_0}\) and dyadic scales
\[
 A\asymp H_A,\quad D\asymp H_D,\quad
 \alpha\asymp E,\quad\beta\asymp F.
\]
There are \(O((\log Y)^4)\) such boxes. Put
\[
 \chi=\frac{A H_D}{D H_A},\quad
 \mu_b=\log_Y\frac{E H_A^2}{u},\quad
 p=\log_Y(H_AH_D),\quad q=\log_Y(EF),\quad z=\log_Yu.
\]
The box and \eqref{eq:old:5_17} imply, uniformly,
\[
 \tau_0-O(1/\log Y)\le\mu_b\le1+O(1/\log Y),\quad
 q_--O(1/\log Y)\le q\le q_++O(1/\log Y),
\]
\[
 p\le z+\mu_b-\frac q2+O(1/\log Y),\qquad 0\le z\le a_0.
\AutoTag{eq:old:5_19}
\]
They also give
\[
 \frac{q_-}{a_0+2/3+e/2}-O(1/\log Y)
 \le \frac qp \le
 \frac{q_+}{2/3-5e/2}+O(1/\log Y).
\AutoTag{eq:old:5_20}
\]
The exact endpoints in \eqref{eq:old:5_20} are
\(19814786/21532607\) and \(20370428/19536965\), both strictly between
\(0.9\) and \(1.1\).

Set
\[
 B_u=\frac98q_+\left(1+a_0-\frac{q_+}{2}\right)
 =\frac{54121135983551}{100000000000000},\qquad
 N_u=\left\lceil Y^{B_u+\xi}\right\rceil .
\AutoTag{eq:old:5_21}
\]
Let \(f(q,z,\mu)=q(1+(z-q/2)/\mu)\). Its derivatives in \(z,\mu,q\)
have signs \(+,+,+\) on \eqref{eq:old:5_19}: for the last sign,
\[
 1+\frac{z-q}{\mu}\ge
 \frac{9351751}{29722179}+O(1/\log Y)>0,
\]
and for the \(\mu\)-derivative use \(q_-/2-a_0>0\). Hence
\[
 \frac98\frac{pq}{\mu_b}\le B_u+O(1/\log Y).
\AutoTag{eq:old:5_22}
\]
Put \(\delta_0=\xi/4\). From \eqref{eq:old:5_19}, \(\mu_b\ge\tau_0-O(1/\log Y)\),
and \(p,q<1+O(1/\log Y)\). Therefore \eqref{eq:old:5_22} gives the direct margin
\[
(B_u+\xi)\mu_b-\frac98(1+\delta_0)pq
\ge \xi\tau_0-\frac{9\xi}{32}-O(1/\log Y)>\frac{\xi}{2}
\]
for all sufficiently large \(Y\).
The thickness in \eqref{eq:old:5_15} satisfies
\(\theta\le2Y^{-\mu_b}\). The strict margin absorbs the factor
\(2\).

Split \(\chi\in(1/2,2)\) into \(N_u\) half-open intervals of length
\(3/(2N_u)\). For \(K_{\rm det}=(r_0+1)(s_0+1)\) source points in one interval,
form the integer matrix whose \(k,(i,j)\)-entry is
\[
 A_k^iD_k^{\,r_0-i}\alpha_k^j\beta_k^{\,s_0-j},
\]
where \(0\le i\le r_0\) and \(0\le j\le s_0\). Choose a fixed
integer \(s_0\) sufficiently large in terms of \(\delta_0\) and the
compact ranges in \eqref{eq:old:5_20}. Put \(r_0=\lceil s_0 q/p\rceil\).
By \eqref{eq:old:5_20}, the possible \(r_0\) form a finite set and
\(r_0/s_0=q/p+O(1/s_0)+O(1/\log Y)\) uniformly. The leading terms
in \eqref{eq:old:5_13} have a uniform positive margin by \eqref{eq:old:5_22}. The ceiling and all
remaining terms in \eqref{eq:old:5_13} are uniformly \(O(1/s_0)+O(1/\log Y)\)
relative to those leading terms. First choose \(s_0\), then the common
threshold for \(Y\). Then these errors use less than half the margin.
Apply the level estimate \eqref{eq:old:5_12}--\eqref{eq:old:5_14} with
\(b_{\rm cell}=B_u+\xi+O(1/\log Y)\) and
\(\mu_{\rm th}=\mu_b+O(1/\log Y)\). It makes every
\(K_{\rm det}\)-row determinant zero for one common sufficiently large \(Y\).
Each entry is an integer of size at most
\[
 (Y+1)^{r_0+(21/20)s_0},
\]
since \(u\le Y^{a_0}<Y^{1/20}\). Hence cofactors give a nonzero
integer polynomial of degree at most \((r_0,s_0)\). Its coefficient
height is \(Y^{C_\xi}\) with fixed \(C_\xi\). This gives the required
integer auxiliary polynomial.

We next count the lattice radii over the moving intervals. Write
\[
\rho=\frac{\alpha/E}{\beta/F},\qquad
\rho=\kappa(1-\theta)\chi^{-2},\qquad
\kappa=\frac{F H_D^2}{E H_A^2}.
\]
Since each normalized coefficient lies in \([1,2)\), actual \(\rho\)
lies in \([1/2,2]\). For each interval, crop its image in \(\rho\) to
\((1/2,2)\). Empty intervals contain no source. The map
\(\chi\mapsto\kappa\chi^{-2}\) has derivative bounded above and below on
\([1/2,2]\), uniformly in every nonempty box. We have
\(1/8\le\kappa\le16\) by the dyadic box inequalities and \eqref{eq:old:5_15}.
Also \(\theta\le2Y^{-\mu_b}<1/N_u\). Hence the retained coefficient directions
have angular width \(O(1/N_u)\). Any angular window of width \(w\)
meets \(O(1+N_uw)\) of them.

For a retained direction with unit vector \(\ell\) and unit normal \(n\),
define
\[
\mathcal A(\alpha,\beta)=
 \left(\ell\cdot(\alpha/E,\beta/F),
       \frac{N_u}{c_0}n\cdot(\alpha/E,\beta/F)\right),
\]
where the fixed \(c_0\) makes the images under \(\mathcal A\)
of the source coefficient vectors lie in the disk of radius
\(5\). This linear map has determinant \(N_u/(c_0EF)\). Its image of the
integer coefficient lattice has a
Euclidean reduced basis \(b_1,b_2\), with angle at least \(60^\circ\).
The inverse images form a unimodular integer basis. For each source pair,
write \((v_1,v_2)\) for its coordinates in this reduced basis. These are
distinct from the original coefficient coordinates \((\alpha,\beta)\). Put
\(T_i=5\sqrt2/|b_i|\). Then \(T_2\le T_1\), every source coefficient
pair has lattice coordinates \(|v_i|\le T_i\), where
\(\mathbf v=(v_1,v_2)\), and
\[
 T_1T_2\gg \frac{EF}{N_u},\qquad T_1^2\gg\frac{EF}{N_u}.
\AutoTag{eq:old:5_23}
\]
The determinant of the lattice and the operator norms of \(\mathcal A\)
and \(\mathcal A^{-1}\) are bounded by fixed powers of \(Y\). Every
nonzero integer preimage has image length at least
\(\|\mathcal A^{-1}\|^{-1}\). Minkowski's theorem bounds the
shortest image vector by a constant times the square root of the
covolume. The reduced-basis angle and determinant then bound the second
basis vector by a fixed power of \(Y\) as well. Apply
\(\mathcal A^{-1}\). Then both unimodular preimage vectors have
polynomially bounded coordinates. Thus the coefficient-lattice change
keeps the auxiliary polynomial's height at \(Y^{O_\xi(1)}\).

We now bound the complete interval sum. Put
\(L_*=\max(1,(EF/N_u)^{1/2})\). Intervals with \(T_1\le2L_*\) contribute
\(O(N_u+(EFN_u)^{1/2})\). For larger \(T_1\), place it in a dyadic range
\((T,2T]\). Let \(\mathbf x_{\min}=(x_1,x_2)\) be the primitive
shortest integer preimage. The definition of \(\mathcal A\) gives
\[
 |(x_1/E,x_2/F)|\ll T^{-1},\qquad
 \operatorname{dist}((x_1/E,x_2/F),\ell)\ll (N_uT)^{-1}.
\]
A nonaxis dyadic annulus of normalized radius \(r\) contains
\(O(r^2EF)\) integer vectors, with
\(r\ge\max(E^{-1},F^{-1})\). A fixed vector is compatible with
\(O(1+(Tr)^{-1})\) interval directions. Sum
\(EF(r^2+r/T)\) over the annuli up to \(r\ll T^{-1}\). This gives
\(O(EF/T^2)\) intervals in this \(T\)-range.
Their contribution to \(\sum(T_1+1)\) is \(O(EF/T)\). The vertical and
horizontal primitive vectors are the two axes. The retained source
directions have coordinate ratio \(\rho\in[1/2,2]\). So they stay a fixed
positive angle from both axes. The corresponding vertical and horizontal
cases cost \(O(F+N_u)\) and \(O(E+N_u)\), respectively. Finally
\(T_1\gg(EF/N_u)^{1/2}\) by \eqref{eq:old:5_23}. Summing the dyadic ranges gives
\[
 \sum_{\text{intervals}}(T_1+1)
 \ll_\eta Y^\eta\{N_u+\sqrt{EFN_u}\log Y+E+F\}.
\AutoTag{eq:old:5_24}
\]

For each interval, factor its auxiliary polynomial over \(\mathbb Q\).
There are \(O_\xi(1)\) factors, each of fixed degree and height \(Y^{O_\xi(1)}\).
Make the unimodular coefficient-lattice change. Then let an irreducible
factor have degree \(a_c\) in the corner pair and \(b_q\) in the
coefficient pair. Write the original coefficients \(\alpha(\mathbf v),\beta(\mathbf v)\)
as homogeneous linear forms in the primitive
lattice coordinates \(\mathbf v=(v_1,v_2)\). The following cases account for
every factor.

\begin{enumerate}
\item If it is not absolutely irreducible, its distinct conjugate
components have no common curve component. By bihomogeneous Bezout they have \(O_\xi(1)\)
rational common points. Each rational point gives at
most four signed primitive pairs.
\item If \(a_c=0\), the factor fixes \(O_\xi(1)\) primitive coefficient
ratios. For each, \eqref{eq:old:5_15} is a square-fibre equation. It has positive
coefficients \(\alpha,\beta\) and nonzero shift \(u\). The coefficients and
shift are at most \((Y+1)^{21/20}\), and \(A\le Y+1\). Theorem~\ref{thm:fibre}
with \(x=2Y\) and fixed height exponent \(A_{\rm ht}=2\) gives \(Y^\eta\)
solutions uniformly.
\item If \(b_q=0\), the factor fixes \(O_\xi(1)\) primitive corner pairs.
The solutions in \((\alpha,\beta)\) to
\(\beta D^2-\alpha A^2=u\) form a progression with step
\((D^2,A^2)\). Its length in the \(\alpha\)-box is at most
\[
 1+\frac{E}{D^2}\ll 1+\frac{uY}{A^2D^2}
 \ll 1+\frac{Y^{-1/3+5e}}u=O(1)
\]
uniformly for \(u\ge1\), by \eqref{eq:old:5_17}--\eqref{eq:old:5_18}.
\item If \(a_c\ge2\) and \(b_q\ge1\), apply Heath--Brown's
Lemma~2. The first pair is the primitive coefficient-lattice pair,
bounded by \(T_1+1\). The second is the primitive corner pair. Its degree
in the second pair is \(a_c\ge2\). Choose the theorem's epsilon after
the fixed polynomial height exponent. Hence the count is
\(O_\eta(Y^\eta(T_1+1))\).
\item If \(a_c=1\) and \(b_q\ge1\), write the factor as
\(A Q(\mathbf v)+D P(\mathbf v)\), with coprime nonzero homogeneous forms
\(P,Q\). The source equation forces \(P(\mathbf v)=-A\zeta\),
\(Q(\mathbf v)=D\zeta\) for a nonzero integer \(\zeta\). The Sylvester
identities put
\(\operatorname{Res}(P,Q)v_i^{2b_q-1}\) in the ideal \((P,Q)\).
Primitivity of \(\mathbf v\) therefore gives
\(\zeta\mid\operatorname{Res}(P,Q)\).
The resultant has fixed polynomial height, so its divisors cost
\(Y^\eta\). Fix a divisor \(\zeta\mid\operatorname{Res}(P,Q)\). Then the
source conic gives the nonzero homogeneous
level equation
\[
\beta(\mathbf v) Q(\mathbf v)^2-\alpha(\mathbf v) P(\mathbf v)^2=u\zeta^2.
\]
The left side is a nonzero binary form of degree \(d_\Phi=2b_q+1\). If it
vanished identically, coprimality of \(P,Q\) would force \(P^2\) to
divide the linear form \(\beta(\mathbf v)\). This is impossible for \(b_q\ge1\).
Denote this form by \(\Phi\). Its level is \(h=u\zeta^2\ne0\). Define
\(\mathcal E=\{v_1\in\mathbb Z:\Phi(v_1,v_2)-h\text{ is identically zero in }v_2\}\).
Writing \(\Phi(v_1,0)=a_\Phi v_1^{d_\Phi}\), every \(v_1\in\mathcal E\) satisfies
\(a_\Phi v_1^{d_\Phi}=h\), so \(|\mathcal E|\le d_\Phi\). If \(a_\Phi=0\), then \(\mathcal E=\varnothing\).
This includes the pure-power case. There every slice is constant, and at most
\(d_\Phi\) slices equal the nonzero target \(h\). There are at most
\(2\lfloor T_1\rfloor+1\) other first coordinates. For each, the nonzero
polynomial in \(v_2\) has degree at most \(d_\Phi\). Hence it has at most \(d_\Phi\)
integer roots. Each exceptional coordinate contributes at most
\(2\lfloor T_2\rfloor+1\) points. Thus
\[
\begin{aligned}
\#\{(v_1,v_2)\in\mathbb Z^2:|v_i|\le T_i,\ \Phi(v_1,v_2)=h\}
&\le d_\Phi(2\lfloor T_1\rfloor+1)+d_\Phi(2\lfloor T_2\rfloor+1)\\
&\le4d_\Phi(T_1+1),\qquad T_2\le T_1.
\end{aligned}
\]
The corners are then forced.
\end{enumerate}

Thus every full interval contributes \(Y^\eta(T_1+1)\). The underfull
intervals cost \(O_\xi(N_u)\). Fix the ordered residual factorizations
\(\alpha=ac,\ \beta=bd\). The inverse formulas \eqref{eq:old:2_2} determine \(B,C\),
then \(m=ABa,n=ACc\). The positivity, integrality, gcd, mask, shell and
first-match filters only discard candidates. The inverse multiplicity is
at most \(\tauDiv(\alpha)\tauDiv(\beta)\ll_\eta Y^\eta\), since
\(\alpha,\beta\le(Y+1)^{21/20}\). The opposite shift is recovered from
\(ut=L\).
Combine this with \eqref{eq:old:5_24}. For one fixed \(u\) and one box the count is
\[
 \ll_\eta Y^\eta\{N_u+\sqrt{EFN_u}\log Y+E+F\}.
\AutoTag{eq:old:5_25}
\]

Split the \(3u\) population counted by \eqref{eq:old:5_25} into
\(\min(A,D)\ge Y^{1/3-\nu}\) and its complement. In the first part,
\eqref{eq:old:5_17} gives
\[
 E+F\ll uY^{1/3+2\nu}.
\AutoTag{eq:old:5_26}
\]
In the second part, suppose \(H_A\le H_D\). The other case is symmetric.
For fixed \(A,\beta\), the congruence
\(\beta D^2\equiv u\pmod{A^2}\) has
\(O(2^{\omega(A)})=Y^\eta\) roots modulo \(A^2\), since
\(\gcd(\beta u,A)=1\). Each root class contributes
\(O(H_D/A^2+1)\) values of \(D\). Also
\(\beta\ll uY/H_D^2\). Summing over \(A,\beta\) gives
\[
    \ll_\eta Y^\eta\left\{\frac{uY}{H_AH_D}
                 +\frac{uYH_A}{H_D^2}\right\}.
\AutoTag{eq:old:5_27}
\]
For odd prime powers the unit-square congruence has at most two roots.
At powers of \(2\) it has at most four. The absolute factor \(4\) is
absorbed in \(O(2^{\omega(A)})\).
By \eqref{eq:old:5_18} and \(H_A\le Y^{1/3-\nu}\), the two terms in \eqref{eq:old:5_27} are at
most
\[
 Y^{1/3+5e/2+o(1)},\qquad
 \frac{Y^{2/3+5e-3\nu+o(1)}}u.
\AutoTag{eq:old:5_28}
\]
The second term sums harmonically over \(u\). The first costs
\(\sum_{u\le Y^{a_0}}1\ll Y^{a_0}\).

Finally \(EF\le M<\kappa_Y^2Y^{q_+}\). Sum \eqref{eq:old:5_25}--\eqref{eq:old:5_28} over
\(u\le Y^{a_0}\), all boxes and the logarithmic factors. Before the final
loss allocation, this gives five exponents and their savings from $2/3$
\[
\begin{array}{@{}ll@{}}
\text{term and exponent} & \text{saving from }2/3\\[8pt]
\text{square-root term: }a_0+(q_++B_u+\xi)/2 &
  \dfrac{5102312049347}{600000000000000}\ (0.00850385\ldots)\\[7pt]
\text{interval term: }a_0+B_u+\xi &
  \dfrac{23206592049347}{300000000000000}\\[7pt]
\text{core axes: }1/3+2\nu+2a_0 & \dfrac{2473}{12000}\\[7pt]
\text{small corner, first: }a_0+1/3+5e/2 & \dfrac{1619393}{6000000}\\[7pt]
\text{small corner, second: }2/3+5e-3\nu & \dfrac{8003}{500000}
\end{array}\AutoTag{eq:old:5_29}
\]
The smallest saving is \(5102312049347/600000000000000>S_0\).
Fix the degree and height constants. Then all fixed-height, divisor,
dyadic and logarithmic losses can be assigned arbitrarily small
exponents. Therefore, for every \(\varepsilon>0\),
\[
 \#3u\ll_\varepsilon
 Y^{2/3-5102312049347/600000000000000+\varepsilon}.
\AutoTag{eq:old:5_30}
\]
The next proposition records this display as the complete first-match row
estimate.
\begin{rowproposition}[\(3u\) row estimate]\label{prop:row-3u}
For every \(\varepsilon>0\), the complete first-match row satisfies
\[
 \#3u\ll_\varepsilon
 Y^{2/3-5102312049347/600000000000000+\varepsilon}.
\]
\end{rowproposition}
\begin{proof}
The full row count is \eqref{eq:old:5_30}. Its five additive terms are bounded in
\eqref{eq:old:5_25}--\eqref{eq:old:5_29}. The least saving is the one displayed there.
\end{proof}

\subsection*{The short-\(t\) row \(3t\)}

In the \(t\)-chart use \(\gamma=ad,\delta=bc\). The exact source equations
are
\[
 \delta C^2-\gamma B^2=t,\qquad
 \frac{\delta}{\gamma}=\left(\frac BC\right)^2(1+\beta_{\rm rem}),\qquad
 \beta_{\rm rem}=\frac{n-m}{m(n+1)}\in(0,1/m).
\AutoTag{eq:old:5_31}
\]
By \eqref{eq:old:5_11}, \(\beta_{\rm rem}\ll Y^{-\tau_0}\). The source heights and the
\(O_R\), \(O_M\), \(Q1\) and \(Q2t\) complements give
\[
 \frac t4Y^{q_-}<BC\le3tY^{c_t},\quad
 \gamma,\delta\ll tY^{g_t},\quad
 \frac{\gamma}{C^2}\le \frac{6Y^{-d_0}}t<1
 \quad\text{for all sufficiently large }Y,
\AutoTag{eq:old:5_32}
\]
where
\[
 c_t=\frac23+l_0=q_+,\quad
 g_t=\frac13+e+l_0+2j_0=\frac{5833463}{15000000},\quad
 d_0=\frac13-2e>0.
\]
For the \(\gamma/C^2\) bound, start from the exact source identity
\[
 n-m=\frac{ZL}{m+n+1}.
\]
Then use
\(BC=(m+n+1)t/L\), \(m+n+1>Y/2\), and
\(L<\kappa_Y Y^{1/3+e}\). Since \(2\kappa_Y\le3\), this gives
\(BC>(t/3)Y^{2/3-e}>(t/4)Y^{q_-}\). For the upper bound use
\(L\ge Y^{1/3-l_0}\) and \(m+n+1\le2Y+1\le3Y\). This gives
\(BC\le3tY^{2/3+l_0}\). Also
\[
 \gamma B^2=\frac{t\,m(n+1)}{n-m},\qquad
 \gamma/C^2=\frac{m(n+1)}{ZBC};
\]
using \(Z>Y^{5/3-e}\) and the lower bound for \(BC\) proves the last
inequality in \eqref{eq:old:5_32}. The bound for \(\delta\) is symmetric. It uses
\(\delta C^2=t\,n(m+1)/(n-m)\) and \(B,C\ge Y^{1/3-j_0}\).

Fix \(t\le Y^{a_0}\) and a dyadic box
\[
 B\asymp H_B,\quad C\asymp H_C,\quad
 \gamma\asymp E,\quad\delta\asymp F.
\]
Put \(\mathcal L=\ln Y\), \(p=\log_Y(8H_BH_C)\),
\(q=\log_Y(128EF)\), \(z=\log_Yt\), and
\[
 \Xi_t=\frac{H_BH_C\sqrt{EF}}{8t},\quad
 \mu_0=\log_Y \Xi_t,\quad
 \ell_1=\ln(8H_BH_C),\quad \ell_2=\ln(128EF).
\]
From \eqref{eq:old:5_11}--\eqref{eq:old:5_32}, with \(\kappa_Y\le3/2\),
\[
 q_-\mathcal L-\ln2<\ell_1\le(c_t+a_0)\mathcal L+\ln24,\quad
 q_-\mathcal L+\ln32<\ell_2<q_+\mathcal L+\ln288,
\]
\[
 \tau_0\mathcal L-\ln256<\ln \Xi_t<(1+e+l_0)\mathcal L,\qquad
 p-z\ge q_--O(1/\mathcal L).
\AutoTag{eq:old:5_33}
\]
Dyadicity gives \(BC/4<H_BH_C\le BC\) and
\(EF\le M<4EF\). With \eqref{eq:old:5_32} these imply
\(8H_BH_C>tY^{q_-}/2\),
\(8H_BH_C\le24tY^{c_t}\),
\(128EF>32Y^{q_-}\), and \(128EF<288Y^{q_+}\).
They also give
\[
 \frac{Y^{\tau_0}}{256}<\Xi_t
 \le\frac{3\kappa_Y}{8}Y^{3q_+/2}
 \le\frac9{16}Y^{1+e+l_0},
\]
which proves the two \(\Xi_t\) bounds in \eqref{eq:old:5_33}.
In particular \(\ell_1,\ell_2,\ln \Xi_t\) lie in
\([\mathcal L/5,2\mathcal L]\) for \(Y\ge2048\).

Set \(\sigma_t=H_B/H_C\) and \(x=(B/C)/\sigma_t\). The source has
\(x\in(1/2,2)\). Partition \(x\in(1/2,2)\) into
\[
 I_k=\left(\frac{k}{N},\frac{k+1}{N}\right].
\]
Equivalently, these are the intervals
\((k\sigma_t/N,(k+1)\sigma_t/N]\) for \(B/C\). Here
\(N/2-1\le k\le2N-1\), with \(N\ge40\) to be fixed below. Let
\(x_4=\lfloor k^2/N\rfloor\) and
\(r_k=x_4\sigma_t^2/N\). For \(x\in I_k\), the width of the interval
and \eqref{eq:old:5_31} give
\[
 |\delta-r_k\gamma|\le18\sigma_t^2E/N,\qquad 0<\gamma<2E.
\AutoTag{eq:old:5_34}
\]
The square of \(B/C\) differs from \(r_k\) by at most
\(5\sigma_t^2/N\). The additional term
\((B/C)^2\beta_{\rm rem}\) is at most \(4\sigma_t^2/m\), and the choice below has
\(N<m\). The endpoint \(k\le2N-1\) follows from \(x<2\). This is the
range used in the \(5\sigma_t^2/N\) bound. A nonempty box also gives
\[
 F/16<\sigma_t^2E<8F.
\AutoTag{eq:old:5_35}
\]
Define the linear map
\[
 \mathcal A_k(y_1,y_2)=
 \left(\frac{N}{\sigma_t^2E}(y_1-r_ky_2),\frac{y_2}{E}\right).
\]
Its image of \(\mathbb Z^2\) has covolume
\(N/(\sigma_t^2E^2)\asymp N/(EF)\). By \eqref{eq:old:5_34}, coefficient pairs
\((\delta,\gamma)\) in the box map into the disk of radius \(19\).
Choose a Euclidean reduced basis \(b_1,b_2\) for this lattice. Let
\(h_1,h_2\) be its unimodular integer preimages. Put
\(T_i=19\sqrt2/|b_i|\). The reduced-basis angle is at least \(60^\circ\).
Hence \(T_2\le T_1\), every source pair has coordinates
\(|v_i|\le T_i\), where \(\mathbf v=(v_1,v_2)\), and \eqref{eq:old:5_35} gives
\[
 T_1T_2
 =\frac{722}{|b_1||b_2|}
 \ge\frac{361\sqrt3}{\operatorname{covol}(\mathcal A_k(\mathbb Z^2))}
 =\frac{361\sqrt3\,\sigma_t^2E^2}{N}
 >\frac{361\sqrt3}{16}\frac{EF}{N}
 >39\frac{EF}{N},\qquad T_1^2\ge T_1T_2.
\AutoTag{eq:old:5_36}
\]
The matrix \(\mathcal A_k\), its inverse, and the covolume have polynomial
height. Every nonzero integer preimage has image length at least
\(\|\mathcal A_k^{-1}\|^{-1}\). Minkowski's theorem bounds the shortest
image vector by a constant times the square root of the covolume. The
reduced-basis angle and determinant then give a polynomial upper bound
for the second vector. Apply \(\mathcal A_k^{-1}\). Then both
unimodular integer preimages have polynomially bounded coordinates.

\begin{samepage}
We now count the moving grid. The integer \(x_4\) lies in
\([N/5,4N]\). Each \(x_4\) comes from at most two \(k\), since
\[
 \sqrt{Nx_4}\le k<\sqrt{N(x_4+1)}
\]
has length at most \(\sqrt5/2<2\) for \(x_4\ge N/5\). Therefore it
contains at most two integers. Take a dyadic range \(T_1\in(T,2T]\). Write
\(\mathcal A_k(y)=b_1\) for the primitive shortest vector
\(y=(y_1,y_2)\). Then
\[
 |y_1-r_ky_2|\ll \sigma_t^2E/(NT),\qquad
 |y_2|\ll E/T.
\AutoTag{eq:old:5_37}
\]
If \(y_2=0\), primitivity gives \(y_1=\pm1\). The contribution
over all \(k\) is then \(O(F+N)\). If \(y_1=0\), then \(y_2=\pm1\).
Since \(x_4\ge N/5\), \eqref{eq:old:5_37} bounds its total contribution by
\(O(E+N)\). For \(y_1y_2\ne0\), \eqref{eq:old:5_37} gives
\(T\ll E\) and \(|y_1|\ll F/T\). Multiply the first inequality in
\eqref{eq:old:5_37} by \(N/\sigma_t^2\). Then, for each fixed \(y_1\), the nonzero
integer \(y_2x_4\) lies in an interval of length \(O(E/T)\). Since
\(E/T\gg1\), there
are \(O(E/T)\) such integer values. Their absolute values satisfy
\(|y_2x_4|\ll NE/T\), a fixed power of \(Y\). So each has \(Y^\eta\) factorizations
as \(y_2x_4\). Each \(x_4\) has at most two \(k\). Thus the number of
intervals in this \(T\)-range is
\(O_\eta(Y^\eta EF/T^2)\). This includes the endpoint cases
\(|y_1|\le C F/T\) and \(E/T\ge c>0\). Their contribution to
\(\sum_k(T_1(k)+1)\) is \(O_\eta(Y^\eta EF/T)\). By \eqref{eq:old:5_36},
\(T\gg(EF/N)^{1/2}\). Sum the dyadic \(T\)-ranges and pay the
\(+1\) terms. This gives
\[
 \sum_k(T_1(k)+1)\ll_\eta
 Y^\eta\{N+\sqrt{EFN}\log Y+E+F\}.
\AutoTag{eq:old:5_38}
\]
\end{samepage}

For each interval \(I_k\), apply the determinant construction \eqref{eq:old:5_12}
to the normalized monomials \(x^i\zeta^j\), where
\(\zeta=(E/F)(\delta/\gamma)=\kappa_t x^2+y\). Scaling rows and columns
back to the integer monomials gives \eqref{eq:old:5_42}. We choose the grid
size \(N\) used above as follows. Put
\[
B_t=\frac34(q_++a_0)
 =\frac{5452607}{10000000},\quad
 \nu_t=B_t+\xi=\frac{5453607}{10000000},\quad
 N_t=\lceil Y^{\nu_t}\rceil .
\AutoTag{eq:old:5_39}
\]
The exact gap
\[
 \tau_0-\nu_t=\frac{2226893}{5000000}>0
\]
shows that \(N_t<m\) for all sufficiently large \(Y\). Here
\(m>Y^{\tau_0}/\sqrt3\) and \(N_t\le2Y^{\nu_t}\). In
\eqref{eq:old:5_34}--\eqref{eq:old:5_38} take \(N=N_t\).
Since \(EF>M/4>Y^{q_-}/4\), we have
\[
\frac{EF}{N_t}\gg Y^{q_- -\nu_t}
=Y^{690793/6000000}.
\]
Thus \eqref{eq:old:5_36} puts every dyadic radius \(T\) above
\(cY^{690793/12000000}>1\) for all sufficiently large \(Y\).
The separately counted \(+1\) per grid interval is already included in
the \(O(N_t)\) term of \eqref{eq:old:5_38}.

The determinant thickness \(\vartheta\) in the normalized coefficient
ratio is comparable to \(\Xi_t^{-1}\). The conic gives
\[
 \frac{\delta}{\gamma}
 =\left(\frac BC\right)^2+\frac{t}{\gamma C^2}.
\]
Normalize the coefficient ratio by \(E/F\). Then the thickness is
\((E/F)t/(\gamma C^2)\asymp t/(F H_C^2)\). From \eqref{eq:old:5_35},
\(F H_C^2\asymp E H_B^2\asymp H_BH_C\sqrt{EF}\). Hence this thickness is
\(\asymp \Xi_t^{-1}\). More precisely, the dyadic bounds
\(E\le\gamma<2E\), \(H_C\le C<2H_C\), and
\(1/16<\sigma_t^2E/F<8\) imply
\[
 \frac{\Xi_t^{-1}}{256}\le\vartheta\le\Xi_t^{-1}.
\]
Thus \(0\le\mu_{\rm th}-\mu_0\le\ln256/\mathcal L\). The normalized
cell width is exactly \(1/N_t\), so
\(\nu_t\le b_{\rm cell}=\log_Y N_t\le\nu_t+\ln2/\mathcal L\).

The box bounds also give
\[
 \frac{q_-}{c_t+a_0}+O(1/\mathcal L)
 \le\frac{\ell_2}{\ell_1}
 \le\frac{q_+}{q_-}+O(1/\mathcal L).
\]
The limiting endpoints are \(9907393/10905214>0.9\) and
\(10185214/9907393<1.1\). So this ratio belongs to \([0.9,1.1]\) for
all sufficiently large \(Y\). Choose a fixed integer \(s_0\)
sufficiently large in terms of the strict margin below and this compact
range. Set
\[
 r_0=\left\lfloor s_0\frac{\ell_2}{\ell_1}\right\rfloor .
\]
The possible \(r_0\) form a finite set. The degree rounding and the
lower-order terms in \eqref{eq:old:5_13} are uniformly \(O(1/s_0)\) relative to the
leading terms. Choose \(s_0\) first. Then enlarge the threshold for \(Y\)
to absorb the \(O(1/\mathcal L)\) errors.

The normalized scale parameters in \eqref{eq:old:5_12} are
\(p=\ell_1/\mathcal L\), \(q=\ell_2/\mathcal L\), and
\(\mu_{\rm th}=p+q/2-z+O(1/\mathcal L)\). Precisely,
\[
 \mu_0=p+q/2-z-\frac{\log(64\sqrt{128})}{\mathcal L}.
\]
The displayed constant and the additional \(O(1/\mathcal L)\) from the
thickness comparison are among the finite errors accounted for below.
The limiting function
\[
 (p,q,z)\longmapsto \frac98\frac{pq}{p+q/2-z}
\]
is increasing in \(p,q,z\) on the source range. The relevant derivative
numerators are \(q(q/2-z)\), \(p(p-z)\) and \(pq\). With
\(D=p+q/2-z\), its three partial derivatives are
\[
\frac98\frac{q(q/2-z)}{D^2},\qquad
\frac98\frac{p(p-z)}{D^2},\qquad
\frac98\frac{pq}{D^2}.
\]
On the stated source domain each has absolute value at most
\((9/8)\cdot4\cdot25=112.5<120\). Here
\[
 q_-/2-a_0=\frac{8467393}{30000000}>0,\qquad
 q_--a_0=\frac{9187393}{15000000}>0.
\]
The bound \(BC\le3tY^{c_t}\) gives the one-sided endpoint estimates
\[
p\le c_t+z+\frac{\ln24}{\mathcal L},\qquad
q\le q_++\frac{\ln288}{\mathcal L}.
\]
First put \(p=c_t+z+O(1/\mathcal L)\). Then the denominator is
\(c_t+q/2+O(1/\mathcal L)\), and \(q/(c_t+q/2)\) increases. Thus the
maximum is \(B_t+O(1/\mathcal L)\), attained at \(q=q_+\), \(z=a_0\).
For \(\delta_0=\xi/4=1/40000\),
\[
 \nu_t-(1+\delta_0)B_t
 =\xi-\delta_0B_t
 =\frac{34547393}{400000000000}>0.
\AutoTag{eq:old:5_40}
\]
Now the uniform error. The upper endpoints in \eqref{eq:old:5_33} change \(p,q\)
from \(c_t+z,q_+\) by at most \(\ln24/\mathcal L,\ln288/\mathcal L\).
The exact normalization of \(\mu_0\) contributes
\(\ln(64\sqrt{128})/\mathcal L\). The thickness comparison contributes
at most \(\ln256/\mathcal L\). Take the source range
\(0\le p,q,z\le2\) and \(p+q/2-z\ge1/5\). There every first derivative of
\(\Phi(p,q,z)=(9/8)pq/(p+q/2-z)\) has absolute value at most \(120\).
Since \(\ln24<4\), \(\ln288<6\), \(\ln(64\sqrt{128})<7\), and
\(\ln256<6\), the total perturbation is at most \(2760/\mathcal L\).
We may therefore take \(C_t=4000\). Hence, whenever
\[
(\xi-\delta_0B_t)\ln Y>2C_t(1+\delta_0),
\AutoTag{eq:old:5_41}
\]
the determinant condition has the strict surplus required in \eqref{eq:old:5_14}.
Take one interval with \(K_{\rm det}=(r_0+1)(s_0+1)\) source rows. Use columns
indexed by \(0\le i\le r_0,\ 0\le j\le s_0\), and multiply each row by
\(C^{r_0}\gamma^{s_0}\). The resulting entry is
\[
 C^{r_0}\gamma^{s_0}(B/C)^i(\delta/\gamma)^j
 =B^iC^{r_0-i}\delta^j\gamma^{s_0-j}\in\mathbb Z.
\AutoTag{eq:old:5_42}
\]
We have \(B,C\le Y+1\) and \(\gamma,\delta\ll Y^{a_0+g_t}\), and the degrees are
fixed. So every matrix entry and cofactor has height \(Y^{O(1)}\). By
\eqref{eq:old:5_12}--\eqref{eq:old:5_14}, every integerized \(K_{\rm det}\times K_{\rm det}\) determinant has absolute
value less than one. Hence it vanishes. If an interval has fewer than
\(K_{\rm det}\) source rows, count those points directly. Otherwise the integer
matrix has a nonzero rational kernel vector. Clearing its denominator
gives a nonzero integral auxiliary polynomial on that interval. It has
bounded bidegree and height \(Y^{O_\xi(1)}\).

In this \(t\)-chart, \(\gcd(\gamma,\delta)=\gcd(B,C)=1\). The
transformed auxiliary factors have height \(Y^{O_\xi(1)}\) by the
fixed-degree factor-height argument above. These are the primitivity and
height hypotheses used in the square-fibre and Heath--Brown counts.
The coefficient-lattice coordinates \(\mathbf v=(v_1,v_2)\) are primitive,
since \((\gamma,\delta)\) is primitive and the basis change is
unimodular. The original coefficients \(\gamma(\mathbf v),\delta(\mathbf v)\)
are homogeneous linear forms in these coordinates.
Factor an auxiliary form over \(\mathbb Q\). Let \(a_c\)
and \(b_q\) be its degrees in the corner and coefficient pairs. There are
five branches. The conjugate-component Bezout argument above handles a
factor that is not absolutely irreducible. If \(a_c=0\),
it fixes \(O_\xi(1)\) primitive coefficient pairs. Then the
square-fibre theorem (Theorem~\ref{thm:fibre} in
Appendix~\ref{app:fibre}) applies to
\(\delta C^2-\gamma B^2=t\), with \(x=2Y\) and fixed height exponent
\(A_{\rm ht}=2\). If \(b_q=0\), it fixes a primitive corner pair
\((B,C)\). Reducing the conic modulo \(C^2\) fixes \(\gamma\). Also
\(0<\gamma<C^2\) by \eqref{eq:old:5_32}. So there is at most one coefficient pair.
If \(a_c\ge2\) and \(b_q\ge1\), apply Heath--Brown's Lemma~2 with the
primitive coefficient pair first. It gives \(O_\eta(Y^\eta(T_1+1))\).
Finally, if \(a_c=1\) and \(b_q\ge1\), write the factor as
\(B Q(\mathbf v)+C P(\mathbf v)\), with coprime nonzero homogeneous forms
\(P,Q\). The conic forces \(P(\mathbf v)=-B\zeta\),
\(Q(\mathbf v)=C\zeta\). The Sylvester
identities and primitivity give \(\zeta\mid\operatorname{Res}(P,Q)\).
There are \(Y^\eta\) choices for \(\zeta\). For each, the conic gives
\[
 \delta(\mathbf v) Q(\mathbf v)^2-\gamma(\mathbf v) P(\mathbf v)^2=t\zeta^2.
\]
  The left side is a nonzero binary form of degree \(2b_q+1\). An identity
  would force \(P^2\) to divide the linear form \(\delta(\mathbf v)\), contrary
  to \(b_q\ge1\). The preceding complete slice count applies with the
  nonzero target \(h=t\zeta^2\). It counts \(O_\eta(T_1+1)\) coefficient
  pairs. Then \(B,C\) are forced. Thus every full interval
  contributes \(Y^\eta(T_1+1)\).

Fix the ordered factorizations \(\gamma=ad,\delta=bc\). The inverse
formulas \eqref{eq:old:2_2} determine
\[
 A=(dB-bC)/t,\qquad D=(cC-aB)/t,
\]
then \(m=ABa,n=ACc\). Also \(u=(cd-ab)/t\). These are forced, not
summed. The factor multiplicity is
\(\tauDiv(\gamma)\tauDiv(\delta)\ll_\eta Y^\eta\). Thus, for fixed \(t\) and
box, \eqref{eq:old:5_38} gives
\[
 \ll_\eta Y^\eta\{N_t+\sqrt{EFN_t}\log Y+E+F\}.
\AutoTag{eq:old:5_43}
\]
The row \(3t\) is contained in this larger population with
\(1\le t\le Y^{a_0}\), \(L\ge Y^{1/3-l_0}\) and \(B,C\ge Y^{1/3-j_0}\).
The restrictions \(t<u\) and \(K_{\rm res}\le Y^{\Delta_{\rm exp}}\) only reduce it.
There are \(O((\log Y)^4)\) boxes. By \eqref{eq:old:5_32}, \(EF\le M\ll Y^{q_+}\)
and \(E+F\ll tY^{g_t}\). Sum \eqref{eq:old:5_43} once over the positive shifts
\(t\le Y^{a_0}\). This gives
\[
 \#3t\ll_\eta Y^\eta\left\{
 Y^{a_0+\nu_t}
 +Y^{a_0+(q_++\nu_t)/2}
 +Y^{g_t+2a_0}\right\}.
\AutoTag{eq:old:5_44}
\]
The three exponents and their savings are
\[
\begin{array}{c|c|c}
\text{term}&\text{exponent}&\text{saving from }2/3\\ \hline
\text{interval count}&5933607/10000000&2199179/30000000\\
\text{square-root count}&39611249/60000000&388751/60000000\ (0.00647918\ldots)\\
\text{coefficient axes}&7273463/15000000&2726537/15000000 .
\end{array}\AutoTag{eq:old:5_45}
\]
The least saving is \(388751/60000000>S_0\). Fix the uniform degree and height constants. Then the fixed loss terms,
divisor factors and logarithms fit within any prescribed
\(\varepsilon>0\). Hence
\[
 \#3t\ll_\varepsilon
 Y^{2/3-388751/60000000+\varepsilon}.
\AutoTag{eq:old:5_46}
\]
\begin{rowproposition}[\(3t\) row estimate]\label{prop:row-3t}
For every \(\varepsilon>0\), the complete first-match row satisfies
\[
 \#3t\ll_\varepsilon
 Y^{2/3-388751/60000000+\varepsilon}.
\]
\end{rowproposition}
\begin{proof}
The complete row is bounded in \eqref{eq:old:5_44}. Its three contributions and
their savings are listed in \eqref{eq:old:5_45}--\eqref{eq:old:5_46}.
\end{proof}
